% !TeX root = ../papers/nc-p3.tex
\section{片側対称な第二補完の階数二曲線の次数}
\label{sec:companion-rank-two-sections}
\begin{proposition}\label{prop:companion-no-binary-quintic}
$(z,\lambda)=(-2,1)$とし、第一行列を
$\left(\begin{smallmatrix}5&2\\3&7\end{smallmatrix}\right)$、
第二行列を$\left(\begin{smallmatrix}1/6&2\\3&-1/4\end{smallmatrix}\right)$
とする。第一関係六平面の階数二以下の軌跡は、素数1009の幾何ファイバーで
次数4と16の二つの整成分を重複度1で持つ。標数零のこの六平面を含む
Grassmann多様体の開集合では、階数二以下の軌跡は次数5の有効一次元
サイクルを含まない。従ってこの辺では両方向とも二変数形式の商が存在しない。
\end{proposition}
\begin{proof}
$\mathbb P(R)=\mathbb P^5$内で関係行列の十六の三次小行列を取る。
$\mathbb F_{1009}$上の三つの超平面切断を
\path{computations/verify_companion_rank_two_sections.py}で計算する。
各切断の$\mathbb P^4$で先頭座標を1とし、残る四変数について
次数5以下の126単項式を並べる。十六の三次式に二次以下の15単項式を掛けた
係数行列の階数は106で、次数5の列は全て主列となる。
従って無限遠超平面に共通零点はなく、全ての高次単項式も四次以下に還元できる。
四行四列の階数二以下の多様体の次数は20で、
$\mathbb P^4$との交叉は固有かつ期待次元0だから切断の長さは20である
（同じ退化軌跡公式をこの四平面に適用できる）。
よって残る20列は座標環の基底である。一座標の乗法の特性多項式は
いずれも平方因子を持たず、その既約因子の次数は順に
\[
 (1,3,16),\qquad(2,2,4,4,8),\qquad(1,4,15)
\]
である。従って各切断は幾何学的に被約で、これらは実際のFrobenius巡回長である。

曲線は期待余次元で、その成分は全て一次元である。これは階数二以下の
行列多様体がCohen--Macaulayで、ここでの線型切断が正則列となることによる。
各成分は超平面に会うため被約切断から一般の重複度は1である。また被約切断の
各点は一つの成分にしか属さない。実際、その点で一つの超平面式を割った
局所環が体なので、曲線の極大イデアルは一元で生成され、正則点である。
Frobeniusが相異なる$r$成分（一つの次数を$e$とする）を推移的に動かすなら、
それらの切断に属する巡回長は$r$の倍数であり、和は$re$である。
どの超平面にも同じ成分の軌道型が適用される。
上の三分割の各々を組に分け、各組の最大公約数の約数$r$を試すと、
共通して残る$(r,e)$のリストは$((1,20))$と$((1,4),(1,16))$だけである。
コードはこの整数の列挙も行う。

薄い二つの関係係数を零にした$\mathbb P^3$では行列がブロック対角となる。
低い二行二列と高い二行二列の各ブロックへの射影は、厚い各色が二座標へ同型に
写るため線型同型である。両行列式が切る二つの滑らかな二次曲面は相異なる。
一致すれば階数二軌跡に曲面が含まれ、既に検算した有限切断に反する。
よってその完全交叉は次数4の有効曲線で、一成分だけの可能性は消える。
特にこの幾何ファイバーはいかなる次数5の有効一次元サイクルも含まない。

固定次数の有効一次元サイクルのChow多様体は固有であり、サイクルの台が
指定した線型空間と階数二以下の閉集合に含まれることは閉条件である。
例えば次数5のサイクルを有限個の成分に分解し、各成分のChow形式で
線型空間と小行列の消滅を課してもよい。従ってサイクルが存在する
六平面の軌跡は閉で、上記の良い還元から有理標本とその開近傍で存在しない。
純関係のない開部分を取ると命題\ref{prop:binary-quotient-rank-two-curve}を
両向きに使える。標本で純関係がないことは、前と同じ二つの角ブロックの行列式と
厚い部分の両射影の単射性から従う。これは非正則性の一つの障害を除く判定である。
\end{proof}
